% Front matter: an associative guide to the notes. New material.
\chapter*{Threads through the notes}
\addcontentsline{toc}{chapter}{Threads through the notes}
\label{THR:ch}

These notes are not meant to be read in a straight line. The same few objects keep turning up in
different clothes, and following them is the most fun way through. Below are some of the threads,
each with its stops. Lectures are numbered $1$--$11$ and the new interludes are lettered A--H.

\begin{center}
\begin{tikzpicture}[x=0.56cm,y=0.52cm]
  \node[font=\small] at (0,1) {1};
  \node[font=\small] at (1,1) {2};
  \node[font=\small] at (2,1) {3};
  \node[font=\small] at (3,1) {4};
  \node[font=\small] at (4,1) {5};
  \node[font=\small] at (5,1) {6};
  \node[font=\small\bfseries,addcol] at (6,1) {A};
  \node[font=\small\bfseries,addcol] at (7,1) {B};
  \node[font=\small\bfseries,addcol] at (8,1) {C};
  \node[font=\small] at (9,1) {7};
  \node[font=\small] at (10,1) {8};
  \node[font=\small] at (11,1) {9};
  \node[font=\small] at (12,1) {10};
  \node[font=\small] at (13,1) {11};
  \node[font=\small\bfseries,addcol] at (14,1) {D};
  \node[font=\small\bfseries,addcol] at (15,1) {E};
  \node[font=\small\bfseries,addcol] at (16,1) {F};
  \node[font=\small\bfseries,addcol] at (17,1) {G};
  \node[font=\small\bfseries,addcol] at (18,1) {H};
  \node[anchor=east,font=\small] at (-0.7,0) {golden ratio};
  \draw[gray!30] (0,0)--(18,0);
  \draw[addcol,thick] (6,0)--(18,0);
  \fill[addcol] (6,0) circle (0.17);
  \fill[addcol] (8,0) circle (0.17);
  \fill[addcol] (13,0) circle (0.17);
  \fill[addcol] (15,0) circle (0.17);
  \fill[addcol] (18,0) circle (0.17);
  \node[anchor=east,font=\small] at (-0.7,-1) {the root system $A_{n-1}$};
  \draw[gray!30] (0,-1)--(18,-1);
  \draw[addcol,thick] (5,-1)--(18,-1);
  \fill[addcol] (5,-1) circle (0.17);
  \fill[addcol] (6,-1) circle (0.17);
  \fill[addcol] (9,-1) circle (0.17);
  \fill[addcol] (10,-1) circle (0.17);
  \fill[addcol] (14,-1) circle (0.17);
  \fill[addcol] (18,-1) circle (0.17);
  \node[anchor=east,font=\small] at (-0.7,-2) {Blichfeldt and Rogers};
  \draw[gray!30] (0,-2)--(18,-2);
  \draw[addcol,thick] (1,-2)--(7,-2);
  \fill[addcol] (1,-2) circle (0.17);
  \fill[addcol] (2,-2) circle (0.17);
  \fill[addcol] (5,-2) circle (0.17);
  \fill[addcol] (7,-2) circle (0.17);
  \node[anchor=east,font=\small] at (-0.7,-3) {$E_8$};
  \draw[gray!30] (0,-3)--(18,-3);
  \draw[addcol,thick] (5,-3)--(18,-3);
  \fill[addcol] (5,-3) circle (0.17);
  \fill[addcol] (6,-3) circle (0.17);
  \fill[addcol] (7,-3) circle (0.17);
  \fill[addcol] (8,-3) circle (0.17);
  \fill[addcol] (15,-3) circle (0.17);
  \fill[addcol] (16,-3) circle (0.17);
  \fill[addcol] (18,-3) circle (0.17);
  \node[anchor=east,font=\small] at (-0.7,-4) {the space of lattices};
  \draw[gray!30] (0,-4)--(18,-4);
  \draw[addcol,thick] (0,-4)--(8,-4);
  \fill[addcol] (0,-4) circle (0.17);
  \fill[addcol] (2,-4) circle (0.17);
  \fill[addcol] (3,-4) circle (0.17);
  \fill[addcol] (7,-4) circle (0.17);
  \fill[addcol] (8,-4) circle (0.17);
  \node[anchor=east,font=\small] at (-0.7,-5) {contacts and critical points};
  \draw[gray!30] (0,-5)--(18,-5);
  \draw[addcol,thick] (12,-5)--(17,-5);
  \fill[addcol] (12,-5) circle (0.17);
  \fill[addcol] (13,-5) circle (0.17);
  \fill[addcol] (14,-5) circle (0.17);
  \fill[addcol] (15,-5) circle (0.17);
  \fill[addcol] (16,-5) circle (0.17);
  \fill[addcol] (17,-5) circle (0.17);
  \node[anchor=east,font=\small] at (-0.7,-6) {entropy};
  \draw[gray!30] (0,-6)--(18,-6);
  \draw[addcol,thick] (2,-6)--(17,-6);
  \fill[addcol] (2,-6) circle (0.17);
  \fill[addcol] (9,-6) circle (0.17);
  \fill[addcol] (13,-6) circle (0.17);
  \fill[addcol] (15,-6) circle (0.17);
  \fill[addcol] (17,-6) circle (0.17);
  \node[anchor=east,font=\small] at (-0.7,-7) {circles and tori};
  \draw[gray!30] (0,-7)--(18,-7);
  \draw[addcol,thick] (6,-7)--(9,-7);
  \fill[addcol] (6,-7) circle (0.17);
  \fill[addcol] (7,-7) circle (0.17);
  \fill[addcol] (8,-7) circle (0.17);
  \fill[addcol] (9,-7) circle (0.17);
  \node[anchor=east,font=\small] at (-0.7,-8) {tetrahedra and icosahedra};
  \draw[gray!30] (0,-8)--(18,-8);
  \draw[addcol,thick] (13,-8)--(18,-8);
  \fill[addcol] (13,-8) circle (0.17);
  \fill[addcol] (15,-8) circle (0.17);
  \fill[addcol] (17,-8) circle (0.17);
  \fill[addcol] (18,-8) circle (0.17);
  \node[anchor=east,font=\small] at (-0.7,-9) {certified computation};
  \draw[gray!30] (0,-9)--(18,-9);
  \draw[addcol,thick] (0,-9)--(16,-9);
  \fill[addcol] (0,-9) circle (0.17);
  \fill[addcol] (13,-9) circle (0.17);
  \fill[addcol] (16,-9) circle (0.17);
  \foreach \x in {5.5,8.5,13.5} \draw[gray,dashed] (\x,1.5)--(\x,-9.5);
\end{tikzpicture}
\end{center}
Chapters are listed in book order (lectures in black, interludes in blue); a dot marks a stop on the thread.


\paragraph{The golden ratio.} Lecture~11 counts the jammed states of a narrow channel with binary words
having no two adjacent $1$s, and Fibonacci numbers appear. The most likely fraction of $1$s is
$(5-\sqrt5)/10$, and the entropy per site is $\log\varphi$. Interlude~E recognizes this as the Shannon
capacity of the ``golden mean'' channel. Interlude~C meets the Fibonacci word again: it is a Sturmian
word of slope $\varphi^{-2}$, the cutting sequence of an irrational line through the square lattice, its gaps obey the three-gap
theorem, and the golden ratio is the worst approximable number (Hurwitz). Folding the $E_8$ diagram
with $\tau=\varphi$ gives the non-crystallographic group $H_4$ and icosahedral quasicrystals
(Interlude~C). Interlude~A sends the binary icosahedral group to the extended $E_8$ diagram (McKay).
The icosahedron itself is not made of regular tetrahedra (Interlude~E).

\paragraph{The root system $A_{n-1}$.} It is built in Lecture~6 as the vectors $e_i-e_j$ in the
hyperplane $\sum x_i=0$. In Lecture~8 the same hyperplanes $x_i=x_j$ cut $\Conf(N,\R)$ into Weyl
chambers. In Interlude~A the roots are the eigenvalues of conjugation on the matrix units
$E_{ij}$; the maximal torus of $\mathrm{SU}(n)$ is $\R^{n-1}$ modulo the $A_{n-1}$ lattice; and its Weyl
alcove is the triangle of gaps between points on a circle from Lecture~7. Complexify the hyperplanes
and divide by $\Sigma_N$, and you get the braid group (Interlude~D).

\paragraph{Blichfeldt.} His 1929 gauge trick bounds the density in Lectures~2 and~3. Rogers' simplex bound (end of Lecture~2) replaces his
spreading of mass by a cutting of Voronoi cells, and is sharp in the plane, where it is Thue's theorem. His 1914 lemma
proves Minkowski's convex body theorem (Interlude~B). The Cohn--Elkies bound is a ``two-sided''
Blichfeldt gauge, with positivity demanded on both sides of the Fourier transform, and with Viazovska's
magic function it settles dimension $8$ (Interlude~B). The gauges for long bodies in Lecture~3 (capsules, cylinders) are the
subject of G.~Fejes T\'oth and W.~Kuperberg, \emph{Blichfeldt's density bound revisited} (Math.\ Ann.\ 1993),
which leads on to the cylinder bounds of A.~Bezdek and W.~Kuperberg (1990) and W.~Kusner (2014).

\paragraph{$E_8$.} It is constructed in Lecture~6 as $D_8\cup(D_8+h)$ and from the Hamming code. In
Interlude~A it is $248=8+240$: Cartan subalgebra plus root spaces, and its density $\pi^4/384$ can be
read off from $\det C=1$. In Interlude~B its theta series is the Eisenstein series $E_4$, so it has
$240\,\sigma_3(m)$ vectors of norm $2m$, and Viazovska's theorem makes it the densest packing in
dimension~$8$. In Interlude~C it folds to $H_4$, and in Interlude~E the Golay code builds its big
sibling, the Leech lattice.

\paragraph{Why root lattices?} They are within a few percent of Rogers' bound in dimensions $2$--$8$
(end of Lecture~6). Their density can be read off from the determinant of the Cartan matrix
(Interlude~A). The Weyl group's irreducibility makes them eutactic and the abundance of roots makes them
perfect, so Voronoi's theorem makes each one a local optimum. Their $O(n^2)$ roots are too few in high
dimension, where Interlude~B's random lattices take over.

\paragraph{The space of lattices.} It appears in Lecture~4 with Mahler's criterion, and in
Interlude~B with Hermite's constant, reduction of quadratic forms and Siegel's mean value theorem. From
the last of these comes the Minkowski--Hlawka lower bound $\dens\ge 2^{-d}\cdot2\zeta(d)$, which was the
lower bound quoted in Lectures~1 and~3. In Interlude~C dynamics on the space takes over: bounded orbits
correspond to badly approximable numbers (Dani), and Margulis's proof of the Oppenheim conjecture.

\paragraph{Contacts, forces and critical points.} The contact graphs and force balance of
Lectures~10--11 make the max-min distance function a ``min-type'' Morse function. Interlude~D turns this
into min-max: the Tammes problem as the top critical value, Betti numbers of configuration spaces as
lower bounds on the number of critical configurations, and kissing numbers. Interlude~E reads the same
force-balance condition as jamming, with Maxwell counting and isostaticity. The slack in the
icosahedral kissing configuration (Interlude~E) is what lets the jitterbug of Lecture~11 move.

\paragraph{Entropy.} Lecture~7 derives the Boltzmann distribution and the Gaussian fluctuations of a
gas in a box. Lecture~11 introduces Edwards' entropy of jammed states and solves hard rods. In
Interlude~E the partition function of hard spheres becomes the volume of a configuration space; the
parking problem has an exact answer; and maximum entropy, Shannon's channel capacity and sphere
packing in high dimensions turn out to be one circle of ideas, with Kabatiansky--Levenshtein
(Lecture~3) as a coding theorem.

\paragraph{Circles and tori.} Points on a circle (Lecture~7) have a triangle of gaps. Interlude~A finds
that triangle as the space of conjugacy classes of $\mathrm{SU}(3)$. Interlude~C cuts the torus along an
irrational line (``shredding the torus''), proves equidistribution, and finds at most three gaps.
Interlude~B counts lattice points in a circle and asks how wrong $\pi r^2$ can be (the Gauss circle
problem).

\paragraph{Tetrahedra and icosahedra.} Twelve balls around one fit at the vertices of an icosahedron
with room to spare. Lecture~11's jitterbug uses that room, and Interlude~E proves it with rotating
tropical rings. Interlude~H sees the same fact in a Coxeter diagram: $\{3,3,5\}$ is spherical, not flat,
so five regular tetrahedra close up around an edge only on $\Sph^3$, in the $600$-cell. In the FCC
lattice, octahedra fill exactly the gap the tetrahedra leave. In the laboratory, tetrahedral dice and
colloidal clusters show what happens in between (Interlude~G).

\paragraph{Certified computation.} Thue's triangle lemma (Lecture~1) is proved again in Interlude~F by
a branch and bound in exact rational arithmetic, using a Lagrange multiplier that is the force balance of
Lectures~10--11. The same chapter treats the case analyses that the Tammes and Kepler problems
need, and the certificates (intervals, sums of squares, formal proof)
that make them trustworthy.
